\documentclass{article}
\usepackage[utf8]{inputenc}
\usepackage[T1]{fontenc} % Mejor manejo de acentos
\usepackage[spanish, es-tabla]{babel}
\usepackage{booktabs} 
\usepackage{caption}  
\usepackage{amsmath}  % Necesario para \text{}

\begin{document}

\begin{table}[ht]
    \centering
    \caption{Auto-calibración de percentiles por dataset y parámetros}
    \label{tab:autocalibracion}
    \small 
    \begin{tabular}{lcccc}
        \toprule
        \textbf{Datasets / Hiperparámetros} & \textbf{Shuttle} & \textbf{Telco Churn} & \textbf{Predict Faults} & \textbf{US Crime} \\
        \midrule
        $PR_{\text{dens}}$ (Radio densidad) & 70 (medio) & 85 (alto) & 85 (alto) & 85 (alto) \\
        Radio físico densidad & 0,40 $\sim$ 981* & 1,43 & 0,582 & 18,2 \\
        $PR_{\text{riesgo}}$ (Radio riesgo) & 30 (muy bajo) & 40 (bajo) & 50 (medio) & 40 (bajo) \\
        Radio físico riesgo & 0,2 $\sim$ 503* & 1 & 0,319 & 6,8 \\
        \addlinespace
        $PR_{\text{entropía}}$ & -- & 70 & -- & 80 \\
        Valor entropía & -- & 0,59 & -- & 0,86 \\
        Criterio pureza & proporción & entropía & proporción & entropía \\
        \addlinespace
        Densidad (umbral densidad) & 25 & 60 & 45 & 45 \\
        Riesgo (umbral riesgo) & 30 & 45 & 55 & 45 \\
        Pureza (umbral proporción) & 40 & -- & 45 & -- \\
        Grado de Isolation Forest & 1 & 0 & 0 & 0 \\
        \bottomrule
        \multicolumn{5}{l}{\textit{Nota.} * Rango por clase.}
    \end{tabular}
\end{table}

\end{document}